\section{Model}

\begin{frame}[allowframebreaks]{模型核心设定与 $(\kappa, \lambda)$ 分离}
两期模型，家庭在 date 0 决定是否接受贷款合约 $(d, p)$，在 date 1 决定是否还款。

\medskip
家庭面临两种不确定性：对商品的价值评价 $\tilde{v}_i$（私人信息）和收入冲击概率 $q_i$（风险类型）。
  
\medskip

\begin{block}{关键}
  传统模型（\textcite{kiyotaki1997credit}）用单一参数刻画抵押品，回收价值和激励效果不可分离。本文将抵押品技术参数化为 $(\kappa, \lambda)$：
  \begin{itemize}
    \item $\kappa$（回收效率）：贷款人在违约后能回收的价值比例。传统担保 $\kappa > 0$；数字担保 $\kappa \approx 0$。
    \item $\lambda$（锁定效力）：违约后借款人只能享受 $(1-\lambda)$ 比例的使用价值。无担保 $\lambda = 0$；完美锁定 $\lambda = 1$。
    \item \colorbox{yellow!40}{数字抵押品的本质:}贷款人几乎收不回任何东西（$\kappa \approx 0$），但借款人的违约代价依然很高（$\lambda$ 接近 1）。
  \end{itemize}
  \end{block}

  \begin{table}[htbp]
  \centering
  \caption{提高回收效率（$\kappa$）与锁定效力（$\lambda$）的比较}
  \label{tab:kappa_vs_lambda}
  \begin{tabular}{l p{5cm} p{5cm}}
    \toprule
    & \textbf{提高 $\kappa$（回收效率）} & \textbf{提高 $\lambda$（锁定效力）} \\
    \midrule
    信贷供给 & 扩大 $\uparrow$ & 扩大 $\uparrow$ \\
    \addlinespace % 增加行间距，让表格更透气
    战略性违约 & \textbf{增加} $\uparrow$ & \textbf{减少} $\downarrow$ \\
    \addlinespace
    机制 & 贷款人回收价值更高 $\rightarrow$ \newline 更愿意放贷 $\rightarrow$ \newline 提高价格 $\rightarrow$ \newline 借款人更多违约 
         & 违约成本更高 $\rightarrow$ \newline 借款人更少违约 $\rightarrow$ \newline 贷款人也因此定更低的价格 \\
    \bottomrule
  \end{tabular}
\end{table}

  \textcite{assuncao2014repossession} 发现巴西简化汽车回收程序后信贷扩张但违约率上升；本文中数字锁定增强后信贷扩张且违约率下降。
\end{frame}

\begin{frame}[allowframebreaks]{两个核心命题：道德风险与逆向选择}

    \begin{questionblock}
      逆向选择和道德风险有什么区别和联系？
    \end{questionblock}

    \begin{block}{道德风险 vs. 逆向选择}
       \begin{table}
  \centering
  \caption{逆向选择与道德风险的比较}
  \label{tab:adverse_selection_vs_moral_hazard}
  \scriptsize
  \begin{tabular}{p{2.2cm} p{4.8cm} p{4.8cm}}
    \toprule
    \textbf{维度} & \textbf{逆向选择（Adverse Selection）} & \textbf{道德风险（Moral Hazard）} \\
    \midrule
    发生时间 & \textbf{事前}：交易达成或契约签订之前 & \textbf{事后}：交易达成或契约签订之后 \\
    不对称类型 & \textbf{隐藏信息}：一方拥有另一方不知道的私人属性或特征 & \textbf{隐藏行动}：一方的行为无法被另一方完全观察或监督 \\
    核心机制 & 劣币驱逐良币，价格机制失去调节作用，导致高质量交易对象退出市场 & 激励扭曲，代理人为自身利益最大化而采取损害委托人利益的行为 \\
    金融学经典场景 & 信贷配给：高利率吓退优质借款人，留下高风险借款人 & 代理成本：经理人在职消费或过度投资 \\
    诺奖代表理论 & 阿克洛夫（Akerlof）的“柠檬市场” & 霍姆斯特罗姆（Holmstrom）的“最优契约理论” \\
    \bottomrule
  \end{tabular}
\end{table}

    \end{block}
    
    \newpage

    \textbf{Date 1 的还款决策}：
    \begin{itemize}
        \item \textbf{还款}：支付 $p$，保留完整的设备服务价值，净效用为 $v_i - p$。
        \item \textbf{违约（被锁定）}：不支付 $p$，设备被远程锁定，只能享受 $(1-\lambda)$ 比例的残余价值，效用为 $(1-\lambda)v_i$。
    \end{itemize}
    $\Rightarrow$ \textbf{还款条件}：$v_i - p \geq (1-\lambda)v_i$，即 $\lambda v_i \geq p$。
    
    \vspace{0.5em}
    
    \begin{block}{Proposition 1：道德风险降低}
        $\lambda$ 越高，\textbf{战略性违约}的概率越低。\\
        \small{机制：不还款意味着失去 $\lambda v_i$ 的服务价值，有能力还款的家庭更倾向于履约。}
    \end{block}
    
    \begin{block}{Proposition 2：逆向选择减轻（正向筛选）}
        存在关于 $\lambda$ 递减的风险阈值 $q_\lambda$，仅 $q_i \leq q_\lambda$ 的家庭接受贷款。\\
        \small{机制：高风险家庭预期被锁定概率更高，担保贷款吸引力降低。}
    \end{block}
    
    \vspace{0.5em}
    
    \begin{itemize}
        \setlength{\itemsep}{4pt}
        \item \textbf{实验映射}：上述两个渠道的分离，直接对应实验中三组比较的识别逻辑。
        \item \textbf{福利的非单调性}：更强的锁定技术不一定提高福利。
        \begin{itemize}
            \item \textit{权衡}：提高 $\lambda$ 会扩大信贷供给，但在非战略性违约（收入冲击）时也会摧毁更多剩余（高风险家庭受害最深）。
            \item \textit{结论}：\textbf{中等程度的 $\lambda$} 可能是福利最优的，这与 PAYGO 实践中相对宽松的锁定政策一致。
        \end{itemize}
    \end{itemize}


    
\end{frame}


% \begin{frame}{两个可检验的异质性预测}
%   \begin{block}{Proposition 5：设备价值的异质性}
%     设备价值越高的家庭，\textbf{接受率更低}，且\textbf{道德风险效应}和\textbf{逆向选择效应}都更大。
%   \end{block}

%   \medskip

%   \begin{block}{Proposition 6：延续价值的异质性}
%     延续价值越高的家庭占比越大，道德风险效应越小。
%   \end{block}

%   \medskip

%   \small
%   这里的“延续价值”指借款人按时还清这笔贷款后，能从与贷款人的持续关系中获得的未来收益。若这部分收益足够高，这些家庭即使没有锁定惩罚，也更倾向于主动还款。
% \end{frame}
